\documentclass[conference]{IEEEtran}

\usepackage{graphicx}
\usepackage{amsmath,amssymb}
\usepackage{float}
\usepackage{booktabs}
\usepackage{multirow}
\usepackage{hyperref}

\newcommand{\LongDate}{%
    \number\day\space
    \ifcase\month
    \or January\or February\or March\or April\or May\or June%
    \or July\or August\or September\or October\or November\or December%
    \fi
    \space \number\year
}

\newcommand{\inputpart}[2]{%
    \vspace{1em}
    \noindent\rule{\linewidth}{0.4pt}
    \section{#1}
    \textit{#2}
    \vspace{0.2em}
}

\begin{document}

\title{
    \textbf{Benchmark Report: 519.lbm\_r} \\
    \normalsize Register Spill Analysis on RISC-V via gem5 \\
    \normalsize Test \textbar{} Train Inputs
}

\author{
    \IEEEauthorblockN{Lütfullah Burak Kaya}
    \IEEEauthorblockA{
        Ozyegin University \\
        burak.kaya.50711@ozu.edu.tr
    }
    \and
    \IEEEauthorblockN{Ismail Akturk}
    \IEEEauthorblockA{
        Ozyegin University \\
        ismail.akturk@ozyegin.edu.tr
    }
}

\twocolumn[
\begin{@twocolumnfalse}
\maketitle
\begin{center}{\small \LongDate}\end{center}
\vspace{1em}
\end{@twocolumnfalse}
]

\begin{abstract}
This report presents the register spill characterisation of
\texttt{519.lbm\_r} across two SPEC CPU2017 input sizes (test, train)
executed on RISC-V (RV64GC) using the gem5 TimingSimpleCPU simulator
with a custom SpillDetector.
Spill rates are \textbf{4.2130\%} (test, 20 time-steps) and
\textbf{4.2229\%} (train, $\sim$300 time-steps).
\texttt{lbm\_r} is the most \emph{stable} benchmark in the study:
the spill rate changes by less than 0.01 percentage points between
input sizes, confirming that the Lattice Boltzmann kernel executes
an identical RISC-V instruction sequence on every time-step regardless
of simulation length.
\end{abstract}

% =========================================================
\section{Benchmark Overview}

\texttt{519.lbm\_r} implements the Lattice Boltzmann Method (LBM) for
3D fluid dynamics simulation.  Each time-step performs collision and
propagation over a $100\!\times\!100\!\times\!130$ grid of fluid cells,
where each cell stores 19 floating-point velocity components (D3Q19
model).

ROI markers surround the main time-step loop in \texttt{src/main.c}:
\begin{itemize}
    \item \texttt{m5\_work\_begin(0,0)} --- before the time-step \texttt{for}-loop
    \item \texttt{m5\_work\_end(0,0)}   --- after  the time-step \texttt{for}-loop
\end{itemize}

\subsection{Input Datasets}

\begin{table}[H]
\centering
\caption{519.lbm\_r Input Dataset Sizes}
\label{tab:inputs}
\begin{tabular}{llr}
\toprule
\textbf{Dataset} & \textbf{Obstacle file} & \textbf{Time-steps} \\
\midrule
test  & \texttt{100\_100\_130\_cf\_a.of} (channel flow)   & 20 \\
train & \texttt{reference.dat}                             & $\sim$300 \\
\bottomrule
\end{tabular}
\end{table}

\subsection{gem5 Configuration (all inputs)}

\begin{itemize}
    \item \textbf{CPU:} TimingSimpleCPU, RISC-V RV64GC
    \item \textbf{Caches:} enabled (default L1 I/D + L2)
    \item \textbf{Memory:} 4\,GB simulated physical RAM
    \item \textbf{Spill mode:} summary (\texttt{spill\_verbose=False})
    \item \textbf{Host:} AMD EPYC 7763 (castEpyc), 128 threads, 251\,GB RAM
\end{itemize}

% ==========================================================
\inputpart{TEST INPUT}{100\_100\_130\_cf\_a.of --- 20 time-steps}
% ==========================================================

\section{Test: Simulation Setup}

\begin{verbatim}
./build/RISCV/gem5.opt
  --outdir=benchmarks/lbm/m5out_test
  configs/deprecated/example/se.py
  --cmd=benchmarks/lbm/lbm_r.riscv
  --options="20 benchmarks/lbm/lbm.in
    0 1 benchmarks/lbm/100_100_130_cf_a.of"
  --cpu-type=TimingSimpleCPU
  --caches --mem-size=4GB
\end{verbatim}

\section{Test: Results}

\subsection{ROI Statistics}

\begin{table}[H]
\centering
\caption{ROI Spill Statistics --- test}
\label{tab:test_roi}
\begin{tabular}{lr}
\toprule
\textbf{Metric} & \textbf{Value} \\
\midrule
ROI instructions    & 6,506,800,041 \\
ROI stores          & 775,731,184 \\
ROI loads           & 1,276,836,306 \\
\midrule
\textbf{ROI spills} & \textbf{274,130,440} \\
\textbf{Spill rate} & \textbf{4.2130\%} \\
Spill / loads       & 21.47\% \\
Spill / stores      & 35.34\% \\
\bottomrule
\end{tabular}
\end{table}

\subsection{Load/Store Breakdown}

\begin{table}[H]
\centering
\caption{ROI Load/Store Breakdown --- test}
\label{tab:test_lsbreakdown}
\begin{tabular}{lrr}
\toprule
\textbf{Category} & \textbf{Count} & \textbf{Share} \\
\midrule
\multicolumn{3}{l}{\textit{Store instructions}} \\
\quad Spill stores     & 274,130,440 & 35.34\% of stores \\
\quad Non-spill stores & 501,600,744 & 64.66\% of stores \\
\quad \textbf{Total stores} & \textbf{775,731,184} & 11.92\% of ROI \\
\midrule
\multicolumn{3}{l}{\textit{Load instructions}} \\
\quad Spill reloads   & 274,130,440 & 21.47\% of loads \\
\quad Non-spill loads & 1,002,705,866 & 78.53\% of loads \\
\quad \textbf{Total loads} & \textbf{1,276,836,306} & 19.62\% of ROI \\
\bottomrule
\end{tabular}
\end{table}

\subsection{Timing}

\begin{table}[H]
\centering
\caption{Simulation Timing --- test}
\label{tab:test_timing}
\begin{tabular}{lr}
\toprule
\textbf{Metric} & \textbf{Value} \\
\midrule
Simulated time           & 37.74\,s \\
Host wall time           & 6,482\,s ($\approx$1.8\,h) \\
Total simulated instr.   & 6,745,885,923 \\
ROI coverage             & 96.5\% \\
\bottomrule
\end{tabular}
\end{table}

\section{Test: Analysis}

The spill rate of \textbf{4.21\%} reflects the LBM kernel's
moderate register pressure: each cell update reads 19 velocity
components, computes collision terms, and writes back 19 updated
values.  The RISC-V ABI's 32 integer and 32 floating-point registers
are insufficient to hold all 19 velocity components plus loop indices
and intermediate sums simultaneously, forcing the compiler to spill a
subset on each cell iteration.

The load-to-store ratio of 1.65:1 ($1{,}276\,\text{M} / 775\,\text{M}$)
is consistent with a stencil computation that reads neighbouring cells
(halo reads) before writing the updated cell.  With Sp/Ld of 21.5\%
and Sp/St of 35.3\%, spill traffic constitutes a meaningful but
sub-dominant fraction of memory operations.

The 3.5\% pre-ROI phase (entry instructions: 239\,M) corresponds to
reading and parsing the obstacle file; this is correctly excluded from
the spill measurement.

% ==========================================================
\inputpart{TRAIN INPUT}{reference.dat --- $\sim$300 time-steps}
% ==========================================================

\section{Train: Simulation Setup}

\begin{verbatim}
./build/RISCV/gem5.opt
  --outdir=benchmarks/lbm/m5out_train
  configs/deprecated/example/se.py
  --cmd=benchmarks/lbm/lbm_r.riscv
  --options="300 benchmarks/lbm/reference.dat
    0 1 benchmarks/lbm/100_100_130_cf_a.of"
  --cpu-type=TimingSimpleCPU
  --caches --mem-size=4GB
\end{verbatim}

\section{Train: Results}

\subsection{ROI Statistics}

\begin{table}[H]
\centering
\caption{ROI Spill Statistics --- train}
\label{tab:train_roi}
\begin{tabular}{lr}
\toprule
\textbf{Metric} & \textbf{Value} \\
\midrule
ROI instructions    & 100,160,472,569 \\
ROI stores          & 11,753,652,015 \\
ROI loads           & 19,574,976,720 \\
\midrule
\textbf{ROI spills} & \textbf{4,229,639,284} \\
\textbf{Spill rate} & \textbf{4.2229\%} \\
Spill / loads       & 21.61\% \\
Spill / stores      & 35.99\% \\
\bottomrule
\end{tabular}
\end{table}

\subsection{Load/Store Breakdown}

\begin{table}[H]
\centering
\caption{ROI Load/Store Breakdown --- train}
\label{tab:train_lsbreakdown}
\begin{tabular}{lrr}
\toprule
\textbf{Category} & \textbf{Count} & \textbf{Share} \\
\midrule
\multicolumn{3}{l}{\textit{Store instructions}} \\
\quad Spill stores     & 4,229,639,284 & 35.99\% of stores \\
\quad Non-spill stores & 7,524,012,731 & 64.01\% of stores \\
\quad \textbf{Total stores} & \textbf{11,753,652,015} & 11.73\% of ROI \\
\midrule
\multicolumn{3}{l}{\textit{Load instructions}} \\
\quad Spill reloads   & 4,229,639,284 & 21.61\% of loads \\
\quad Non-spill loads & 15,345,337,436 & 78.39\% of loads \\
\quad \textbf{Total loads} & \textbf{19,574,976,720} & 19.54\% of ROI \\
\bottomrule
\end{tabular}
\end{table}

\subsection{Timing}

\begin{table}[H]
\centering
\caption{Simulation Timing --- train}
\label{tab:train_timing}
\begin{tabular}{lr}
\toprule
\textbf{Metric} & \textbf{Value} \\
\midrule
Simulated time           & 551.01\,s \\
Host wall time           & 164,226\,s ($\approx$45.6\,h) \\
Total simulated instr.   & 100,400,367,313 \\
ROI coverage             & 99.8\% \\
\bottomrule
\end{tabular}
\end{table}

\section{Train: Analysis}

The train spill rate of \textbf{4.22\%} is virtually identical to the
test result (difference: $+0.01$ pp).  This confirms that the LBM
kernel is \emph{perfectly regular}: each of the $\sim$300 time-steps
executes the same RISC-V instruction sequence over the same grid,
producing the same register pressure.  No warm-up, branch
misprediction, or data-dependent code-path variation affects the
spill count.

The train ROI covers 99.8\% of total simulated instructions, indicating
that virtually all computation is within the timed region.  The
load-to-store ratio remains 1.66:1 and the Sp/Ld/Sp/St ratios shift
by less than 1 pp compared to test, further confirming the kernel's
computational regularity.

% ==========================================================
\section*{Cross-Input Summary}

The following metrics are used:
\begin{itemize}
    \item \textbf{Rate} $= \text{ROI spills} / \text{ROI instructions}$
    \item \textbf{Sp/Ld} $= \text{ROI spills} / \text{ROI loads}$
    \item \textbf{Sp/St} $= \text{ROI spills} / \text{ROI stores}$
\end{itemize}
Spill counts are abbreviated: B\,=\,$10^9$, M\,=\,$10^6$.

\begin{table}[H]
\centering
\caption{519.lbm\_r --- All Runs Across Input Sizes}
\label{tab:summary}
\resizebox{\columnwidth}{!}{%
\begin{tabular}{lrrrr}
\toprule
\textbf{Input / Steps} & \textbf{Spills}
    & \textbf{Rate} & \textbf{Sp/Ld} & \textbf{Sp/St} \\
\midrule
\multicolumn{5}{l}{\textit{Test}} \\
\quad 20 steps (cf\_a.of) & 274\,M & 4.2130\% & 21.47\% & 35.34\% \\
\midrule
\multicolumn{5}{l}{\textit{Train}} \\
\quad $\sim$300 steps (ref.dat) & 4.23\,B & 4.2229\% & 21.61\% & 35.99\% \\
\bottomrule
\end{tabular}%
}
\end{table}

\texttt{lbm\_r} is unique among all 17 SPEC CPU2017 benchmarks evaluated:
the spill rate is \emph{input-invariant} to four significant figures.
The $\Delta\text{Rate} = +0.0099$ pp between test and train is within
measurement noise, confirming that LBM register pressure is determined
solely by the kernel's register allocation, not by the number of
time-steps or obstacle configuration.  The reference input (3000
time-steps) is expected to produce an identical rate when completed.

\end{document}